\section*{Chapitre \ref{ch:nw:networking-for-trading} --- Réseaux pour le trading}

\begin{solution}{exo:nw:networking-for-trading:1}
$8 \times (256 + 20) / 25 = \qty{88.3}{\nano\second}$. Une trame de 64~octets occupe 84~octets de temps de ligne, soit 672~bits~:
$10^{10}/672 \approx 14.88$ millions de trames par seconde.
\end{solution}

\begin{solution}{exo:nw:networking-for-trading:2}
Les 23~bits de poids faible de 233.54.12.111 sont 54.12.111 (54 est inférieur à 128, donc rien n'est perdu), ce qui donne l'adresse
\texttt{01:00:5e:36:0c:6f}. Tout groupe
ayant les mêmes 23~bits de poids faible la partage, par exemple 224.54.12.111 ou 233.182.12.111 ($182 = 54 + 128$).
\end{solution}

\begin{solution}{exo:nw:networking-for-trading:3}
$2 \times 1.1 \times 50.1 = \qty{110.2}{\giga\bit\per\second}$ au pic de 10~millisecondes~; $110.2/25 = 4.4$, donc cinq
liens de \qty{25}{\giga\bit\per\second} (ou deux à 100).
\end{solution}

\begin{solution}{exo:nw:networking-for-trading:4}
Avec l'\omterm{def:nw:networking-for-trading:nagle}{algorithme de Nagle} activé, l'annulation est retenue tant que le nouvel ordre n'est pas acquitté~; la plateforme l'acquitte avec son message suivant
vers la société (l'acceptation de l'ordre) ou lorsque sa minuterie d'acquittement différé expire, ce que TCP autorise à durer jusqu'à
une demi-seconde. L'annulation est retardée, jamais réordonnée~: TCP livre les octets dans l'ordre où ils ont été écrits. \texttt{TCP\_NODELAY} supprime
l'attente.
\end{solution}

\begin{solution}{exo:nw:networking-for-trading:5}
L'excédent vaut $4 \times 6 - 10 = \qty{14}{\giga\bit\per\second}$ pendant \qty{50}{\micro\second}~: 700\,000~bits, soit 87\,500~octets. Le
dernier octet attend $87\,500 \times 8 / 10 = 70\,000$~ns, soit \qty{70}{\micro\second}.
\end{solution}

\begin{solution}{exo:nw:networking-for-trading:6}
Janvier 2024~: $80 / 42.4 = 1.9$~; juillet 2026~: $350 / 63.9 = 5.5$. Le chiffre à la seconde a augmenté de moitié, la milliseconde la plus chargée a plus que
quadruplé~: les rafales sont devenues plus brutales. Une liaison dimensionnée en 2024 pour la milliseconde la plus chargée de cette année (80~millions de messages par seconde)
écoule moins du quart de la milliseconde la plus chargée de juillet 2026 (350~millions).
\end{solution}

\begin{solution}{exo:nw:networking-for-trading:7}
Environ 0.32~Mo lorsque 10~\% des paquets appartiennent à des grappes communes, 1.10~Mo à 30~\% et 1.98~Mo à 50~\% (moyennes sur vingt graines). La
mémoire tampon croît à peu près proportionnellement à cette part~: la taille des grappes communes, et donc le retard que laisse chaque événement, lui est
proportionnelle, alors que les rafales indépendantes seules nécessitent 26~Ko.
\end{solution}

\begin{solution}{exo:nw:networking-for-trading:8}
Un pic d'une minute est une moyenne sur soixante mille millisecondes~; la perte survient pendant les quelques millisecondes, voire moins, durant lesquelles le
port reçoit plus de \qty{10}{\giga\bit\per\second}, ce que le compteur ne peut pas montrer. Il faut vérifier les compteurs de rejets en sortie
du commutateur pour ce port et mesurer le trafic à la résolution de la milliseconde et de la microseconde (les captures du chapitre~5) avant
d'incriminer le serveur.
\end{solution}

\begin{solution}{pb:nw:networking-for-trading:1}
\begin{enumerate}
\item $8 \times 0.6 = \qty{4.8}{\giga\bit\per\second}$~: 48~\% du port.
\item Chaque trame occupe 220~octets de temps de ligne, soit 1\,760~bits~: $4.8 \times 10^9 / 1\,760 \approx 2.73$ millions de trames par seconde.
\item $8 \times 2.5 = \qty{20}{\giga\bit\per\second}$, soit le double du \omterm{def:nw:networking-for-trading:ser}{débit de ligne}~: un taux d'utilisation de 2.
\item Une milliseconde d'excédent dans une minute déplace la moyenne sur une minute de moins de 0.002~\%.
\item $(20 - 10) \times 10^9 \times 10^{-3} / 8 = 1.25$~Mo.
\item $1\,250\,000 / 220 \approx 5\,682$ trames.
\item $8 \times 1.25 \times 10^6 / 10^{10} = \qty{1}{\milli\second}$.
\item Elle se vide à $10 - 4.8 = \qty{5.2}{\giga\bit\per\second}$~: $10^7 / 5.2 \times 10^9 \approx \qty{1.92}{\milli\second}$.
\item 1.25~Mo.
\item $6\,000\,000 / 16 = 375\,000$~octets.
\item $1\,250\,000 - 375\,000 = 875\,000$~octets, soit environ 3\,977 trames.
\item La rafale présente $20 \times 10^9 \times 10^{-3} / 1\,760 \approx 11\,364$ trames~: 35~\% sont rejetées.
\item Oui~: 6~Mo est supérieur à 1.25~Mo, si le port peut emprunter la totalité du bloc.
\item \textbf{Résultat nommé.} Une perte nulle exige 1.25~Mo de mémoire tampon pour ce port~; avec son seizième du bloc partagé, soit 375~Ko, cet
événement rejette environ 3\,977 trames, soit 35~\% de la rafale.
\item $6 \times 10^6 \times 8 / 10^{10} = \qty{4.8}{\milli\second}$ de la même rafale.
\item Non~: 20 est inférieur à 25.
\item Non~: les copies arrivent en même temps, sur des ports du même bloc, si bien que chacune obtient la même part et que les deux lignes perdent le même
type de trames au même moment~; les lignes doivent passer par des commutateurs ou des blocs différents.
\item Le compteur de rejets en sortie du commutateur pour le port (par file)~; le compteur de trous du feed handler, avec les trous comblés par aucune des deux lignes.
\item Les trous sur les deux lignes sont envoyés au serveur de retransmission s'il les possède dans sa fenêtre et ses limites de débit, sinon au
prochain instantané~; dans les deux cas, le carnet d'ordres reste périmé entre-temps.
\item Parce que l'événement de marché qui fait rafaler un flux fait rafaler tous les flux, et qu'une mémoire tampon partagée n'est vide que lorsque les
rafales ne se produisent pas.
\end{enumerate}
\end{solution}

\begin{solution}{iq:nw:networking-for-trading:1}
Le temps nécessaire pour cadencer une trame sur le câble~: $8(\ell_{\mathrm f}+20)/R_{\mathrm L}$, soit 67.2~ns pour la plus petite trame à 10, 26.9 à
25. Il est payé à l'émetteur, puis à chaque saut en stockage et réexpédition, indépendamment de la distance, ce qui en fait une part importante
d'un chemin court.
\iqlookfor{la formule avec les 20~octets de surcharge, et le fait qu'elle se répète à chaque saut.}
\end{solution}

\begin{solution}{iq:nw:networking-for-trading:2}
Les données de marché vont d'un vers plusieurs et ne doivent ralentir pour personne~: le multicast envoie une seule copie que le réseau réplique, et UDP
n'a aucun contrôle de flux~; les pertes sont gérées par des lignes redondantes, la retransmission et les instantanés. Les ordres vont d'un vers un et doivent arriver
exactement une fois et dans l'ordre, ce que garantit TCP~; la couche session s'appuie dessus.
\iqlookfor{un vers plusieurs contre un vers un, et l'endroit où la récupération des pertes se situe dans chaque cas.}
\end{solution}

\begin{solution}{iq:nw:networking-for-trading:3}
Nagle retient les nouvelles petites données tant que les données antérieures ne sont pas acquittées~; le récepteur peut retarder son acquittement jusqu'à une minuterie. Ensemble,
un second petit ordre peut attendre l'expiration de la minuterie. Il faut activer \texttt{TCP\_NODELAY}, écrire chaque message en un seul appel, garder la connexion active,
et acquitter sans délai là où on contrôle le côté récepteur.
\iqlookfor{l'interaction des deux mécanismes, et pas seulement le nom de l'option.}
\end{solution}

\begin{solution}{iq:nw:networking-for-trading:4}
En amont du serveur~: le port du commutateur qui y mène (rejets en sortie pendant les \omterm{def:nw:networking-for-trading:burst}{micro-rafales}), le chemin de la couche 1 ou du commutateur des deux lignes
si elles le partagent, le passage à la plateforme. Sur le serveur~: les pertes de la mémoire tampon de réception du noyau, que la carte ne comptabilise pas. Il faut regarder
les compteurs de rejets par port et une capture horodatée au point de prélèvement.
\iqlookfor{une cause commune aux deux lignes et une recherche qui suit le chemin du paquet.}
\end{solution}

\begin{solution}{iq:nw:networking-for-trading:5}
Il lit les rapports d'adhésion IGMP et n'envoie chaque groupe qu'aux ports qui s'y sont abonnés. Sans lui, chaque trame multicast est
inondée sur tous les ports~: les liens et les serveurs transportent et rejettent un trafic qu'ils n'ont pas demandé, et les \omterm{def:nw:networking-for-trading:burst}{micro-rafales} tombent sur des ports qui
auraient dû rester calmes.
\iqlookfor{l'inondation comme défaillance, et son coût en bande passante et en processeur.}
\end{solution}

\begin{solution}{iq:nw:networking-for-trading:6}
À partir des rafales corrélées, et non de la moyenne~: il faut mesurer l'arrivée conjointe des flux à la résolution de la microseconde les jours chargés, calculer le
volume en attente de la borne fluide ou rejouer la capture à travers un modèle de file, retenir un quantile élevé du volume de pointe, et le comparer
avec la mémoire tampon que le port peut réellement obtenir lorsque ses voisins sont eux aussi chargés. Si elle ne suffit pas, augmenter le débit du port ou répartir
les flux.
\iqlookfor{la corrélation entre flux, les mémoires tampons partagées et une rafale mesurée plutôt que supposée.}
\end{solution}
